\begin{table*}[t]
  \centering
  \small
  \begin{tabularx}{\textwidth}{@{}
    >{\raggedright\arraybackslash}p{.17\textwidth}
    >{\raggedright\arraybackslash}p{.22\textwidth}
    >{\raggedright\arraybackslash}p{.18\textwidth}
    >{\raggedright\arraybackslash}X@{}}
    \toprule
    Condition & Construction & Information category & Held-out behaviour \\
    \midrule
    Selected structural
      & Greedy/backward cover of declared candidates
      & Inductive; development cells only
      & Frozen rules may activate on new combinations. \\
    Factorised--development-frozen
      & Expert factor rules retained only with development item support
      & Inductive; expert hypothesis family
      & Supported reusable factors may transfer; unsupported factors are absent. \\
    Full-cell--development-frozen
      & One KC for each exact development cell
      & Inductive exact-cell control
      & Held-out cells have zero active KCs and remain in evaluation. \\
    Predefined factorised
      & Versioned expert rules for reusable marked factors
      & Expert prior over the full domain
      & Factor rules activate wherever their predicates match. \\
    Factorised plus interactions
      & Predefined factors plus named structural interactions
      & Expert interaction prior
      & Both reusable factors and matching interactions can activate. \\
    Full-cell--all
      & One KC for every exact cell in the inventory
      & Transductive inventory control
      & Every held-out cell is known as its own KC; no component reuse. \\
    \bottomrule
  \end{tabularx}
  \caption{Ontology conditions and their prior-information categories.  The
  transductive control is descriptive and must not be compared with inductive
  methods as if it had the same information.}
  \label{tab:ontology-conditions}
\end{table*}
